\documentclass[a4paper, serif, 10pt]{moderncv}

%% moderncv
% color and style
\moderncvcolor{black}
\moderncvstyle{banking}

%% page layout/margins
\usepackage[top=2.5cm, bottom=2.5cm, left=1.5cm, right=1.5cm]{geometry}

\nopagenumbers{}

%% Babel config for Dutch language
% ref:
% - https://latex3.github.io/babel/guides/locale-dutch.html
\usepackage[dutch]{babel}

%% fontspec
\usepackage{fontspec}

\IfFontExistsTF{Linux Libertine}{
  \setmainfont[Ligatures={TeX, Common}, Numbers=OldStyle]{Linux Libertine}
}{}
\IfFontExistsTF{Linux Libertine O}{
  \setmainfont[Ligatures={TeX, Common}, Numbers=OldStyle]{Linux Libertine O}
}{}

%% CTeX
% CJK support, used to show CJK characters in the resume
%
% - heading=false: disable ctex heading
% - punct=kaiming: use kaiming punctuations styles for CJK
% - scheme=plain: use plain scheme, do not override \normalsize font size
% - space=auto: space settings for CJK characters
%
% CJK fonts are resolved per locale language, e.g. Noto Serif/Sans CJK JP
% for Japanese resumes, with fallback to the next region if not installed
%
% ref:
% - http://ctan.mirrorcatalogs.com/language/chinese/ctex/ctex.pdf
\usepackage[UTF8, heading=false, punct=kaiming, scheme=plain, space=auto]{ctex}

\IfFontExistsTF{Noto Serif CJK SC}{
  \setCJKmainfont{Noto Serif CJK SC}
}{}
\IfFontExistsTF{Noto Sans CJK SC}{
  \setCJKsansfont{Noto Sans CJK SC}
}{}

%% line spacing
\usepackage{setspace}
\setstretch{1.125}

%% URL styling - use normal text instead of monospace
\urlstyle{same}

%% Patch \httplink and \httpslink to handle full URLs
% The original moderncv macros always prepend http:// or https:// to URLs.
% This causes issues when the URL already has a protocol (e.g., https://example.com)
% resulting in malformed URLs like https://https://example.com.
% This patch checks if the URL already contains "://" and uses it directly if so.
% Note: We use \str_set:Nx (x-expansion) to expand macros like \@homepage first.
\ExplSyntaxOn
\str_new:N \g_yamlresume_url_str

\RenewDocumentCommand{\httpslink}{O{}m}{
  \str_set:Nx \g_yamlresume_url_str {#2}
  \str_if_in:NnTF \g_yamlresume_url_str {://}
    { \url{#2} }
    { \url{https://#2} }
}

\RenewDocumentCommand{\httplink}{O{}m}{
  \str_set:Nx \g_yamlresume_url_str {#2}
  \str_if_in:NnTF \g_yamlresume_url_str {://}
    { \url{#2} }
    { \url{http://#2} }
}
\ExplSyntaxOff

\name{Daan van den Berg}{}
\title{Werktuigbouwkundig ingenieur met focus op precisietechniek en duurzame productontwikkeling}
\phone[mobile]{+31 6 12345678}
\email{daan.vandenberg@example.nl}
\homepage{https://www.daanvandenberg.nl}

\address{Oudegracht 123, Utrecht, Utrecht, Nederland, 3512 AB}{}{}

\extrainfo{{\small \faLinkedin}\ \href{https://www.linkedin.com/in/daanvandenberg}{@daanvandenberg} {} {} {} • {} {} {} 
{\small \faGithub}\ \href{https://github.com/daanvandenberg}{@daanvandenberg}}

\begin{document}

\maketitle

\section{Basis Informatie}

\cvline{}{\begin{itemize}
\item Ervaren werktuigbouwkundig ingenieur met een masterdiploma van de TU Delft en meer dan vijf jaar ervaring in de hightech- en maakindustrie.
\item Sterk in het ontwerpen en valideren van complexe mechanische systemen, van concept tot productie, met oog voor kosten, kwaliteit en duurzaamheid.
\item Bedreven in CAD (SolidWorks, CATIA V5, Siemens NX), FEM-analyse (ANSYS) en data-gedreven optimalisatie met MATLAB en Python.
\item Teamspeler met uitstekende communicatieve vaardigheden en ervaring in multidisciplinaire projecten volgens Agile/Scrum en Lean-principes.
\end{itemize}}
\section{Opleidingen}

\cventry{sep 2017–jul 2019}
        {Master, Werktuigbouwkunde, Score: 8.2/10}
        {Technische Universiteit Delft}
        {\url{https://www.tudelft.nl}}
        {}
        {\begin{itemize}
\item Afgestudeerd op het ontwerp en de validatie van een compacte warmtewisselaar voor industriële koelsystemen.
\item Specialisatie in precisietechniek, materiaalkunde en geavanceerde simulatiemethoden.
\item Projecten uitgevoerd in samenwerking met ASML, Demcon en VDL.
\end{itemize}
\textbf{Cursussen}: Mechanica van Materialen, Numerieke Stromingsleer, Precisietechniek, Productietechniek, Systeemontwerp, Regeltechniek}

\cventry{sep 2013–jul 2017}
        {Bachelor, Werktuigbouwkunde, Score: 8.0}
        {Hogeschool Utrecht}
        {\url{https://www.hu.nl}}
        {}
        {\begin{itemize}
\item Breed opgeleide werktuigbouwkundige met praktijkervaring in ontwerp, productie en onderhoud.
\item Minor in Duurzame Energie en stage bij Fokker Technologies.
\item Actief lid van de studievereniging en deelnemer aan nationale ontwerpwedstrijden.
\end{itemize}
\textbf{Cursussen}: Sterkteleer, Warmteleer, CAD-Technieken, Productieprocessen, Elektrotechniek, Projectmatig Werken}
\section{Werkervaring}

\cventry{jan 2022–Heden}
        {Senior Werktuigbouwkundig Ingenieur}
        {ASML}
        {\url{https://www.asml.com}}
        {}
        {\begin{itemize}
\item Verantwoordelijk voor het mechanisch ontwerp en de integratie van precisiemodules voor geavanceerde lithografiesystemen.
\item Uitvoeren van FEM-analyses en tolerantieanalyses om de stabiliteit en nauwkeurigheid van systemen te waarborgen.
\item Samenwerken met elektronica-, software- en productie-engineeringteams in een multidisciplinaire omgeving.
\item Begeleiden van junior ingenieurs en bijdragen aan continue verbetering van ontwerpprocessen.
\end{itemize}
\textbf{Trefwoorden}: Precisietechniek, FEM-analyse, Cleanroom, Systeemtechniek}

\cventry{sep 2019–dec 2021}
        {Werktuigbouwkundig Ingenieur}
        {VDL Enabling Technologies Group}
        {\url{https://www.vdlgroep.com}}
        {}
        {\begin{itemize}
\item Ontworpen van machineonderdelen en assemblages voor de halfgeleider- en medische industrie.
\item Opgesteld van technische specificaties, tolerantieanalyses en productierapporten.
\item Begeleid van prototyping en serieproductie bij interne en externe leveranciers.
\item Verantwoordelijk voor kwaliteitsverbeteringen en kostenreductie binnen lopende projecten.
\end{itemize}
\textbf{Trefwoorden}: Machinebouw, GD\&T, Prototyping, Kwaliteitsbeheer}

\cventry{feb 2016–jul 2016}
        {Stagiair Werktuigbouwkunde}
        {Fokker Technologies}
        {\url{https://www.fokker.com}}
        {}
        {\begin{itemize}
\item Ondersteund bij het ontwerp en de analyse van lichte vliegtuigcomponenten.
\item Uitgevoerd van materiaalonderzoek en het opstellen van testprotocollen.
\item Gerapporteerd over resultaten aan het engineeringteam en de stagebegeleider.
\end{itemize}
\textbf{Trefwoorden}: Luchtvaart, Materiaalkunde, Testen, Rapportage}
\section{Talen}

\cvline{Nederlands}{Moedertaal of tweetalige vaardigheid}
\cvline{Engels}{Volledige professionele vaardigheid \hfill \textbf{Trefwoorden}: TOEFL 110, IELTS 7.5}
\cvline{Duits}{Beperkte werkvaardigheid \hfill \textbf{Trefwoorden}: Goethe B1}
\section{Vaardigheden}

\cvline{Werktuigbouwkundig Ontwerp}{Expert \hfill \textbf{Trefwoorden}: SolidWorks, CATIA V5, Siemens NX, GD\&T, DFM}
\cvline{Simulatie en Analyse}{Geavanceerd \hfill \textbf{Trefwoorden}: ANSYS, FEM, MATLAB, Python, CFD}
\cvline{Productie en Prototyping}{Geavanceerd \hfill \textbf{Trefwoorden}: 3D-printen, CNC-frezen, Lasersnijden, Assemblage}
\cvline{Projectmanagement}{Intermediair \hfill \textbf{Trefwoorden}: Agile, Scrum, Lean Six Sigma, IPMA}
\section{Onderscheidingen}

\cventry{jul 2019}
        {Technische Universiteit Delft}
        {Beste Afstudeerproject Werktuigbouwkunde}
        {}
        {}
        {Toegekend voor het ontwerp van een compacte, energiezuinige warmtewisselaar met een verbeterde warmteoverdracht van 18\%.}

\cventry{nov 2021}
        {VDL Groep}
        {Innovatieprijs VDL Groep}
        {}
        {}
        {Prijs voor een kostenbesparend ontwerp van een module voor de halfgeleiderindustrie, met een reductie van assemblagekosten van 12\%.}
\section{Certificaten}

\cventry{mrt 2020}
        {Dassault Systèmes}
        {Certified SolidWorks Professional (CSWP)}
        {\url{https://www.solidworks.com/certification}}
        {}
        {}

\cventry{jun 2021}
        {ASQ}
        {Lean Six Sigma Green Belt}
        {\url{https://asq.org/cert/lean-six-sigma-green-belt}}
        {}
        {}

\cventry{sep 2022}
        {IPMA Nederland}
        {Projectmanagement IPMA-D}
        {\url{https://www.ipma.nl}}
        {}
        {}
\section{Publicaties}

\cventry{mei 2021}
        {Optimalisatie van warmtewisselaars voor compacte industriële systemen}
        {TVVL Magazine}
        {\url{https://www.tvvl.nl}}
        {}
        {\begin{itemize}
\item Beschrijft een methode om warmtewisselaars te optimaliseren met behulp van CFD en experimenteel onderzoek.
\item Toont een verbetering van de warmteoverdracht en een reductie van het drukverlies in industriële koelsystemen.
\item Gepresenteerd op de TVVL Innovatiedag 2021.
\end{itemize}}
\section{Projecten}

\cventry{sep 2018–jun 2019}
        {Ontwerp en validatie van een compacte warmtewisselaar voor industriële koelsystemen.}
        {Compacte Warmtewisselaar}
        {\url{https://www.tudelft.nl}}
        {}
        {\begin{itemize}
\item Ontworpen in SolidWorks en geanalyseerd met ANSYS Fluent.
\item Verbeterde warmteoverdracht met 18\% en verlaagde drukval met 10\%.
\item Prototype getest in het laboratorium van de TU Delft.
\end{itemize}
\textbf{Trefwoorden}: Warmteoverdracht, CFD, SolidWorks, Prototyping}

\cventry{feb 2018–aug 2018}
        {Ontwikkeling van een modulaire robotarm voor flexibele productietoepassingen.}
        {Modulaire Robotarm}
        {\url{https://www.hu.nl}}
        {}
        {\begin{itemize}
\item Ontworpen van een robotarm met verwisselbare grijpers en een bereik van 1,2 meter.
\item Geïntegreerd van sensoren en een eenvoudige besturing met Python.
\item Gepresenteerd op de Techniekbeurs Utrecht 2018.
\end{itemize}
\textbf{Trefwoorden}: Robotica, Modulair Ontwerp, Python, Sensoren}
\section{Interesses}

\cvline{Techniek en Innovatie}{3D-printen, Robotica, Duurzame Energie}
\cvline{Sport en Buitenleven}{Wielrennen, Wandelen, Zeilen}
\section{Vrijwilligerswerk}

\cventry{sep 2020–Heden}
        {Techniekmentor}
        {Stichting Techniekpromotie}
        {\url{https://www.techniekpromotie.nl}}
        {}
        {\begin{itemize}
\item Begeleiden van middelbare scholieren bij techniekprojecten en het ontwerpen van eenvoudige machines.
\item Organiseren van workshops over 3D-printen en duurzame energie.
\item Bijdragen aan het vergroten van interesse in techniek en werktuigbouwkunde onder jongeren.
\end{itemize}}
    
\end{document}